% Generated draft: requires mathematical and source review.
\documentclass[11pt]{article}
\usepackage[margin=25mm]{geometry}
\usepackage{fontspec}
\setmainfont{Latin Modern Roman}
\usepackage{amsmath,amssymb,mathtools,unicode-math}
\setmathfont{Latin Modern Math}
\usepackage{enumitem,needspace,xurl,hyperref}
\setlength{\parindent}{0pt}
\setlength{\parskip}{5pt}
\setlength{\emergencystretch}{3em}
\newcommand{\N}{\mathbb{N}}
\newcommand{\Z}{\mathbb{Z}}
\newcommand{\Q}{\mathbb{Q}}
\newcommand{\R}{\mathbb{R}}
\newcommand{\C}{\mathbb{C}}
\newcommand{\F}{\mathbb{F}}
\newcommand{\A}{\mathbb{A}}
\newcommand{\PP}{\mathbb{P}}
\newcommand{\E}{\mathbb{E}}
\newcommand{\eps}{\varepsilon}
\newcommand{\dd}{\,\mathrm{d}}
\begin{document}
\section*{UnsolvedMath ID: 30000585}
\textbf{Source problem code:} catalogue entry
\subsection{Definitions and mathematical statement}
The source asks for the order of the force-induced compact-to-expanded phase transition in a partially directed self-attracting polymer model on \(\Z^2\). It names a generating function \(H(x,y,z)\) recording contacts, span in the force direction and length, but supplies neither an explicit function nor a complete definition of admissible walks.

Using the indicated order of variables as a notation convention, let \(\mathcal W_n\) be the intended length-\(n\) polymer configurations, \(c(\omega)\) their contact count and \(s(\omega)\) their force-direction span. For positive contact and span weights define
\[
 Z_n(x,y)=\sum_{\omega\in\mathcal W_n}x^{c(\omega)}y^{s(\omega)},\qquad
 H(x,y,z)=\sum_{n\geq0}Z_n(x,y)z^n.
\]
This spells out the generating-function role conditionally; it does not specify the step set, self-avoidance/contact rules, anchoring, boundary constraints or physical fugacity convention omitted by the record. The span is not silently replaced by an endpoint displacement.

With \(f=\log y\), a standard dimensionless force variable, the relevant free energy would be
\[
 \kappa(x,f)=\lim_{n\to\infty}\frac1n\log Z_n(x,e^f),
\]
provided the thermodynamic limit exists. Fixing an admissible attractive regime \(x\), the question asks for the singularity as \(f\) crosses its compact-to-expanded threshold \(f_c(x)\). In a conventional force ensemble the conjugate observable is extension or span per monomer: at finite \(n\),
\[
 \frac{\partial}{\partial f}\frac1n\log Z_n(x,e^f)
 =\frac1n\mathbb E_{n,x,f}[s].
\]
A jump in the limiting force derivative/extension describes a first-order transition; continuous extension with a singular higher derivative or other critical behavior requires the corresponding order or exponent convention. No interchange of differentiation and the thermodynamic limit, existence of a threshold, or actual jump/continuity is asserted here.

The source does not specify the force direction, contact/temperature range, physical identification of \(x,y\), definition of compact versus expanded phases, or whether special endpoints have different orders. Thus neither a universal first- nor second-order answer can be attached to this underspecified model, and no thermal transition is substituted for the force-induced one.

The literature triage checked 2026-08-21 reports no authoritative resolution of the extracted claim located in its search. The open label and notation-only verification note are retained as dated metadata, not independently checked current results.
\subsection{Short English statement}
Determine how the polymer free energy becomes singular when force drives the intended compact polymer into an expanded phase, in particular whether the conjugate extension jumps. The contact, span and length generating variables are retained, while the walk rules, parameter regime and precise transition-order convention remain missing.
\subsection{Sources}
\begin{itemize}
\item Local UnsolvedMath registry, permanent ID 30000585, source code catalogue entry, ``Force-Induced Phase Transitions in Self-Attracting Polymers,'' including original and cleaned model descriptors, generating-function role and verification metadata.
\item \emph{Spatial Random Processes and Statistical Mechanics (2007)}, Oberwolfach Reports, pp.\ 2447--2504, DOI 10.4171/owr/2006/41: \url{https://doi.org/10.4171/owr/2006/41}. This is also the sole supplied literature reference. The citation uses 2007, while the DOI and proposed-year metadata use 2006.
\item Literature triage checked 2026-08-21; local record updated 2026-09-16. No external source has been visited or independently checked here.
\end{itemize}
\end{document}
